\newpage{\ }
\thispagestyle{empty}
\chapter*{Resumen}

\noindent El presente Trabajo Fin de Grado consiste en el desarrollo de una aplicación móvil orientada al análisis del fútbol profesional, centrada en la Primera División de la Liga Española. El objetivo principal del proyecto es ofrecer una herramienta accesible, visual e interactiva que permita consultar información deportiva de forma sencilla, combinando datos estadísticos, análisis avanzados y funcionalidades predictivas.

La aplicación permite al usuario explorar información sobre temporadas, equipos, jugadores y partidos, ofreciendo una visión completa de la competición. A través de distintas pantallas y \textit{dashboards}, el usuario puede consultar clasificaciones, calendarios, rankings, estadísticas individuales y colectivas, trayectorias de equipos, rendimiento de jugadores y análisis específicos de partidos. De esta forma, la información deportiva se presenta de manera clara y organizada, evitando que el usuario tenga que interpretar grandes volúmenes de datos sin procesar.

Uno de los aspectos más relevantes de la aplicación es la incorporación de una sección específica de análisis avanzado. En ella se muestran resultados generados mediante técnicas de minería de datos y aprendizaje automático, como predicciones de partidos, simulaciones de clasificación, estados de forma de equipos y jugadores, jugadores similares, probables goleadores y recomendaciones de fichajes. Estas funcionalidades permiten al usuario ir más allá de la consulta estadística tradicional y obtener una interpretación más profunda del rendimiento deportivo.

Además, la aplicación incluye un asistente conversacional denominado Agente de oro, que permite realizar preguntas en lenguaje natural sobre la información almacenada en el sistema. Gracias a esta funcionalidad, el usuario puede consultar datos concretos sin necesidad de navegar manualmente por todas las pantallas, haciendo que la experiencia sea más directa e intuitiva.

El proyecto busca acercar al usuario herramientas de análisis que habitualmente están reservadas a entornos profesionales, como clubes, analistas deportivos u organizaciones especializadas. Para ello, se ha diseñado una aplicación que combina información histórica, visualización de datos, predicciones y consultas inteligentes en una única plataforma. En conjunto, el trabajo demuestra cómo el uso de datos puede ayudar a comprender mejor el fútbol, facilitando el acceso a información útil tanto para aficionados como para perfiles más técnicos interesados en el análisis deportivo.

\vspace{.5cm}

\noindent \textbf{Palabras clave:} fútbol, LaLiga, análisis deportivo, minería de datos, predicción de partidos, \textit{dashboards}, inteligencia artificial, aplicación móvil, \textit{business intelligence}, aprendizaje automático.